\chapter{Multiplication and the Conway induction}

\section{The recursive product}

\begin{definition}[Product and product-option expression]
  \label{def:game-mul}
  \lean{Game.mul,Game.mulOpt4,mem_mul_left,mem_mul_right}
  \leanok
  \uses{def:game,def:birthday,lem:birthday-options,def:game-inductions,def:game-add,def:game-neg}
  Put
  \[
    \mulopt(x',y;x,y')=x'y+xy'-x'y'.
  \]
  For \(x=\{X^L\mid X^R\}\) and \(y=\{Y^L\mid Y^R\}\), multiplication is the recursive cut
  \[
    xy=
    \left\{
      \begin{array}{ll}
        \mulopt(x^L,y;x,y^L),&\mulopt(x^R,y;x,y^R)
      \end{array}
      \ \middle|\
      \begin{array}{ll}
        \mulopt(x^L,y;x,y^R),&\mulopt(x^R,y;x,y^L)
      \end{array}
    \right\},
  \]
  with all choices of the displayed options.  By
  Lemma~\ref{lem:birthday-options}, each selected option has smaller birthday
  than its parent.  Hence the two-game induction relation of
  Definition~\ref{def:game-inductions} decreases in every recursive call.
\end{definition}

\begin{theorem}[Zero, one, and commutativity for game multiplication]
  \label{thm:game-mul-laws}
  \lean{Game.mul_zero_eq,Game.zero_mul_eq,Game.mul_zero,Game.zero_mul,Game.mul_comm,Game.mul_one_eq,Game.mul_one,Game.one_mul}
  \leanok
  For all games \(a,b\),
  \[
    a0\GameEq0,\qquad 0a\GameEq0,\qquad ab\GameEq ba,
    \qquad a1\GameEq a,\qquad 1a\GameEq a.
  \]
  In fact, multiplication by zero on either side and multiplication on the right by one
  satisfy literal game-tree equalities.
\end{theorem}
\begin{proof}
  \leanok
  \uses{def:game,lem:birthday-options,def:game-inductions,thm:game-order-calculus,thm:option-extensionality,thm:game-add-laws,thm:game-add-order,def:game-neg,thm:game-neg-order,def:game-mul}
  In \(a0\), the second factor has no options.  Expanding the recursive cut in
  Definition~\ref{def:game-mul} therefore gives two empty option lists, so
  \(a0=0\) as literal game trees.
  The same direct expansion of \(0a\) has no option in any of the four
  families, because the first factor has no options; hence \(0a=0\) literally
  as well.

  For commutativity, use the two-game well-founded induction of
  Definition~\ref{def:game-inductions}; every recursive pair is smaller by
  Lemma~\ref{lem:birthday-options}.  Split an option according to the four
  families in Definition~\ref{def:game-mul}.  For a same-side example, the
  \(LL\) option of
  \(ab\) determined by \((a^L,b^L)\) is
  \[
    \mulopt(a^L,b;a,b^L)=a^Lb+ab^L-a^Lb^L.
  \]
  The \(LL\) option of \(ba\) determined by \((b^L,a^L)\) is
  \[
    \mulopt(b^L,a;b,a^L)=b^La+ba^L-b^La^L.
  \]
  The induction hypotheses for the three smaller pairs
  \((a^L,b)\), \((a,b^L)\), and \((a^L,b^L)\) identify, respectively,
  \(a^Lb\GameEq ba^L\), \(ab^L\GameEq b^La\), and
  \(a^Lb^L\GameEq b^La^L\).  Additive commutativity is
  Theorem~\ref{thm:game-add-laws}, compatibility of addition with
  \(\GameEq\) is part of Theorem~\ref{thm:game-add-order}, and compatibility
  of negation with \(\GameEq\) is part of
  Theorem~\ref{thm:game-neg-order}.  These three results therefore make the
  two displayed options equivalent.  The \(RR\) case is the identical
  calculation with
  \(a^L,b^L\) replaced by \(a^R,b^R\):
  \[
    \mulopt(a^R,b;a,b^R)\GameEq\mulopt(b^R,a;b,a^R).
  \]

  A mixed-side option changes family when the factors are swapped.  The
  \(LR\) option
  \[
    \mulopt(a^L,b;a,b^R)=a^Lb+ab^R-a^Lb^R
  \]
  is matched not with another \(LR\) option, but with the \(RL\) option
  \[
    \mulopt(b^R,a;b,a^L)=b^Ra+ba^L-b^Ra^L
  \]
  of \(ba\).  The induction hypotheses for \((a^L,b)\), \((a,b^R)\), and
  \((a^L,b^R)\) identify its three terms after the same additive
  rearrangement.  Dually,
  \[
    \mulopt(a^R,b;a,b^L)\GameEq\mulopt(b^L,a;b,a^R),
  \]
  so an \(RL\) option of \(ab\) matches an \(LR\) option of \(ba\).
  Repeating these matches from \(ba\) back to \(ab\), for both the left and
  right option lists, supplies the four hypotheses of
  Theorem~\ref{thm:option-extensionality}; hence \(ab\GameEq ba\).  The two
  literal zero identities already prove both displayed zero laws.

  Finally, prove the literal equality \(a1=a\) by the one-game induction of
  Definition~\ref{def:game-inductions}.  By Definition~\ref{def:game},
  \(1=\{0\mid\}\).  Thus an \(LL\) option of \(a1\) is
  \[
    \mulopt(a^L,1;a,0)=a^L1+a0-a^L0,
  \]
  and an \(RL\) option is the same expression with \(a^R\) in place of
  \(a^L\); the other two families are empty.  The induction hypothesis gives
  \(a^L1=a^L\) and \(a^R1=a^R\), while the zero-product identity established
  above, the additive zero laws in Theorem~\ref{thm:game-add-laws}, and the
  identity \(-0=0\) in Definition~\ref{def:game-neg} reduce the displayed
  expressions to \(a^L\) and \(a^R\).  Hence the left and right option lists of \(a1\) are
  literally those of \(a\).  The commutativity conclusion already proved then
  gives \(1a\GameEq a\).
\end{proof}
\begin{theorem}[Distributivity on games]
  \label{thm:game-distrib}
  \lean{Game.mul_distrib_tri,Game.mul_distrib,Game.mul_distrib_right}
  \leanok
  For all games \(a,b,c\),
  \[
    a(b+c)\GameEq ab+ac,
    \qquad
    (a+b)c\GameEq ac+bc.
  \]
\end{theorem}
\begin{proof}
  \leanok
  \uses{def:game-inductions,lem:birthday-options,thm:game-order-calculus,thm:option-extensionality,def:game-add,thm:game-add-order,thm:game-neg-order,def:game-quotient,thm:game-quotient-ordered-add-group,def:game-mul,thm:game-mul-laws}
  Prove the first equivalence by the three-game well-founded induction of
  Definition~\ref{def:game-inductions}.  Lemma~\ref{lem:birthday-options}
  shows that every recursive triple is smaller.  For an option of
  \(a(b+c)\), Definition~\ref{def:game-mul} first distinguishes whether the
  changed coordinate is \(a\) or \(b+c\); in the second case,
  Definition~\ref{def:game-add} distinguishes whether the changed coordinate
  is \(b\) or \(c\).  The matching option of \(ab+ac\) changes the same
  coordinate in the corresponding summand.  Conversely, an option of the sum
  changes either \(ab\) or \(ac\), and Definition~\ref{def:game-mul}
  determines the corresponding option on the left.  This gives four
  left/right matching assertions, two in each direction.

  For example, let \(a^L\) and \(b^L\) be left options.  The corresponding
  \(LL\) option on the left is
  \[
    L=\mulopt(a^L,b+c;a,b^L+c).
  \]
  The matching left option on the right is
  \[
    L'=\mulopt(a^L,b;a,b^L)+ac.
  \]
  The three smaller induction hypotheses give
  \[
    \begin{aligned}
      a^L(b+c)&\GameEq a^Lb+a^Lc,\\
      a(b^L+c)&\GameEq ab^L+ac,\\
      a^L(b^L+c)&\GameEq a^Lb^L+a^Lc.
    \end{aligned}
  \]
  Substitution in the definition of \(\mulopt\) gives
  \[
    L\GameEq
      (a^Lb+a^Lc)+(ab^L+ac)-(a^Lb^L+a^Lc)
      \GameEq \mulopt(a^L,b;a,b^L)+ac=L'.
  \]
  The cancellation in the second step takes place in the auxiliary quotient of
  Definition~\ref{def:game-quotient}, with its ordered additive group structure from
  Theorem~\ref{thm:game-quotient-ordered-add-group}.  Every remaining
  membership branch uses one of the same two algebraic templates, according
  as the changed summand is \(b\) or \(c\); reverse-direction branches use
  symmetry of the resulting equivalence.  Left and right options are placed
  as prescribed by Definition~\ref{def:game-mul}.  Compatibility with sums and negatives
  and additive-group normalization follow from
  Theorem~\ref{thm:game-quotient-ordered-add-group}, which puts each pair of signed sums
  in the same form.  By the
  quotient relation in
  Definition~\ref{def:game-quotient}, equality of their classes is precisely
  game equivalence.  The four matches now satisfy
  Theorem~\ref{thm:option-extensionality}, proving
  \(a(b+c)\GameEq ab+ac\).

  Finally,
  \[
    (a+b)c\GameEq c(a+b)\GameEq ca+cb\GameEq ac+bc.
  \]
  The first equivalence uses product commutativity from
  Theorem~\ref{thm:game-mul-laws}; the middle equivalence is the
  left-distributive result just proved; and the last uses product
  commutativity inside the two summands together with addition congruence from
  Theorem~\ref{thm:game-add-order}.
\end{proof}

\begin{lemma}[Multiplication and negation]
  \label{lem:game-mul-neg}
  \lean{Game.mul_neg,Game.neg_mul}
  \leanok
  For all games \(u,v\),
  \[
    u(-v)\GameEq-(uv),
    \qquad
    (-u)v\GameEq-(uv).
  \]
\end{lemma}
\begin{proof}
  \leanok
  \uses{def:game-inductions,lem:birthday-options,def:game-neg,def:game-mul,
    thm:game-order-calculus,thm:option-extensionality,thm:game-neg-order,
    def:game-quotient,thm:game-quotient-ordered-add-group,thm:game-mul-laws}
  Prove the first equivalence by the two-game induction of
  Definition~\ref{def:game-inductions}.  Let \(u'\) be either a left or a
  right option of \(u\), and let \(v'\) be either a left or a right option of
  \(v\).  The three smaller induction hypotheses are
  \[
    u'(-v)\GameEq-(u'v),\qquad
    u(-v')\GameEq-(uv'),\qquad
    u'(-v')\GameEq-(u'v').
  \]
  In the auxiliary quotient, these hypotheses give
  \[
    \begin{aligned}
      [\mulopt(u',-v;u,-v')]
        &=[u'(-v)]+[u(-v')]-[u'(-v')]\\
        &=-[u'v]-[uv']+[u'v']\\
        &=-\bigl([u'v]+[uv']-[u'v']\bigr)\\
        &=[-\mulopt(u',v;u,v')].
    \end{aligned}
  \]
  Equality of the first and last classes means
  \[
    \mulopt(u',-v;u,-v')
      \GameEq-\mulopt(u',v;u,v').
    \tag{*}
  \]

  Definition~\ref{def:game-neg} gives
  \(({-v})^L=-v^R\) and \(({-v})^R=-v^L\).  Consequently, applying \((*)\)
  to the indicated choices gives the two left-option matches
  \[
    \begin{aligned}
      \mulopt(u^L,-v;u,-v^R)
        &\GameEq-\mulopt(u^L,v;u,v^R),\\
      \mulopt(u^R,-v;u,-v^L)
        &\GameEq-\mulopt(u^R,v;u,v^L),
    \end{aligned}
  \]
  and the two right-option matches
  \[
    \begin{aligned}
      \mulopt(u^L,-v;u,-v^L)
        &\GameEq-\mulopt(u^L,v;u,v^L),\\
      \mulopt(u^R,-v;u,-v^R)
        &\GameEq-\mulopt(u^R,v;u,v^R).
    \end{aligned}
  \]
  The expressions on the right in the first display are precisely the
  negatives of the two right options of \(uv\), hence are left options of
  \(-(uv)\); those in the second display are the negatives of the two left
  options of \(uv\), hence are right options of \(-(uv)\).  Reading the same
  matches in reverse supplies the converse directions.  Thus
  Theorem~\ref{thm:option-extensionality} proves
  \(u(-v)\GameEq-(uv)\).

  Finally,
  \[
    (-u)v
      \GameEq v(-u)
      \GameEq-(vu)
      \GameEq-(uv).
  \]
  The first and last steps use product commutativity from
  Theorem~\ref{thm:game-mul-laws}, with negation congruence from
  Theorem~\ref{thm:game-neg-order} in the last step; the middle step is the
  first equivalence.
\end{proof}

\section{A simultaneous well-founded proof}

\begin{definition}[Dershowitz--Manna goal order]
  \label{def:dm-goals}
  \lean{Goal,μA,μB,μGoal,GoalLT,wfGoal}
  \leanok
  \uses{def:birthday,lem:birthday-options}
  Three kinds of induction goal are used:
  \[
    A(x,y),\qquad B(x_1,x_2,y),\qquad C(x_1,x_2,y).
  \]
  Their measures are the multisets
  \[
    \mu_A(x,y)=\{\bd(x),\bd(y)\},\qquad
    \mu_B(x_1,x_2,y)=\mu_C(x_1,x_2,y)
      =\{\bd(x_1),\bd(x_2),\bd(y)\}.
  \]
  The goal order is the pullback of the Dershowitz--Manna multiset extension
  \(<_{\rm DM}\) of \(<\) on \(\mathbb N\) along these measures.  For example,
  if \(a'<a\) and \(a_1,a_2<a\), then
  \[
    \{a',b\}<_{\rm DM}\{a,b\},\qquad
    \{a',b,c\}<_{\rm DM}\{a,b,c\},
  \]
  \[
    \{a_1,a_2,b\}<_{\rm DM}\{a,b\},\qquad
    \{a,b\}<_{\rm DM}\{a,b,c\}.
  \]
  The first two comparisons replace one birthday by a smaller birthday, the
  third replaces one birthday by two smaller birthdays, and the fourth deletes
  one birthday.  Lemma~\ref{lem:birthday-options}, transitivity, and permutation
  invariance of multisets reduce the recursive dependencies in the
  \(A\)-, \(B\)-, and \(C\)-cases to these examples.  Since \(<\) on
  \(\mathbb N\) is well founded, so is the resulting goal order.
\end{definition}

\begin{definition}[Cross inequalities]
  \label{def:cross-conditions}
  \lean{CLeft,CRight,MulCrossLe,MulCrossLt}
  \leanok
  \uses{def:game-order,def:game-add,def:game-mul}
  For \(x_1\GameLt x_2\), the two forms needed to compare product options are
  \[
    \begin{aligned}
      C_L(x_1,x_2;y,y^L)&:\quad
        x_1y+x_2y^L\GameLt x_1y^L+x_2y,\\
      C_R(x_1,x_2;y,y^R)&:\quad
        x_1y^R+x_2y\GameLt x_1y+x_2y^R.
    \end{aligned}
  \]
  For arbitrary games \(x_1,x_2,y_1,y_2\), define the weak and strict
  four-product relations by
  \[
    \begin{aligned}
      \operatorname{Cross}_{\GameLe}(x_1,x_2;y_1,y_2)&:\quad
        x_1y_2+x_2y_1\GameLe x_1y_1+x_2y_2,\\
      \operatorname{Cross}_{\GameLt}(x_1,x_2;y_1,y_2)&:\quad
        x_1y_2+x_2y_1\GameLt x_1y_1+x_2y_2.
    \end{aligned}
  \]
  Thus \(C_L(x_1,x_2;y,y^L)\) is
  \(\operatorname{Cross}_{\GameLt}(x_1,x_2;y^L,y)\), while
  \(C_R(x_1,x_2;y,y^R)\) is
  \(\operatorname{Cross}_{\GameLt}(x_1,x_2;y,y^R)\).
\end{definition}

\begin{theorem}[Cross-inequality calculus]
  \label{thm:cross-calculus}
  \lean{MulCrossLt.swap,MulCrossLt.trans,MulCrossLt.congr_right,MulCrossLt.congr_left,mulOpt4_xslot_mono_le_of_cross,mulOpt4_xslot_antitone_le_of_cross,mulOpt4_yslot_mono_le_of_cross,mulOpt4_yslot_antitone_le_of_cross,mulOpt4_xslot_mono_lt_of_cross,mulOpt4_xslot_antitone_lt_of_cross,mulOpt4_yslot_mono_lt_of_cross,mulOpt4_yslot_antitone_lt_of_cross}
  \leanok
  Let \(\mathrel{\bowtie}\) denote either \(\GameLe\) or \(\GameLt\).  If
  \(\operatorname{Cross}_{\bowtie}(a,c;d,e)\) holds, then
  \[
    \begin{array}{ll}
      \text{(a)}&
        \mulopt(a,e;b,d)\mathrel{\bowtie}\mulopt(c,e;b,d),\\
      \text{(b)}&
        \mulopt(c,d;b,e)\mathrel{\bowtie}\mulopt(a,d;b,e),\\
      \text{(c)}&
        \mulopt(a,b;c,d)\mathrel{\bowtie}\mulopt(a,b;c,e),\\
      \text{(d)}&
        \mulopt(c,b;a,e)\mathrel{\bowtie}\mulopt(c,b;a,d).
    \end{array}
  \]
  For the strict relation, the transpose, composition, and two endpoint-replacement rules are
  \[
    \begin{aligned}
      \text{(e)}\quad&
      \operatorname{Cross}_{\GameLt}(a,c;d,e)
        \Longrightarrow \operatorname{Cross}_{\GameLt}(d,e;a,c),\\
      \text{(f)}\quad&
      \operatorname{Cross}_{\GameLt}(a,b;d,e)\ \wedge
      \operatorname{Cross}_{\GameLt}(b,c;d,e)
        \Longrightarrow \operatorname{Cross}_{\GameLt}(a,c;d,e),\\
      \text{(g)}\quad&
      \operatorname{Cross}_{\GameLt}(a,b;d,e)\ \wedge
      bd\GameEq cd\ \wedge\ be\GameEq ce
        \Longrightarrow \operatorname{Cross}_{\GameLt}(a,c;d,e),\\
      \text{(h)}\quad&
      ad\GameEq bd\ \wedge\ ae\GameEq be\ \wedge
      \operatorname{Cross}_{\GameLt}(b,c;d,e)
        \Longrightarrow \operatorname{Cross}_{\GameLt}(a,c;d,e).
    \end{aligned}
  \]
\end{theorem}
\begin{proof}
  \leanok
  \uses{def:game-add,thm:game-add-laws,def:game-mul,thm:game-mul-laws,def:game-quotient,thm:game-quotient-ordered-add-group,def:cross-conditions}
  Work in the quotient \(\mathcal G/{\GameEq}\) of
  Definition~\ref{def:game-quotient}, equipped with the ordered additive group structure of
  Theorem~\ref{thm:game-quotient-ordered-add-group}.  The hypothesis is
  \[
    [ae]+[cd]\mathrel{\bowtie}[ad]+[ce].
  \]
  Adding the same class to both sides and rearranging the additive group gives,
  respectively,
  \[
    \begin{aligned}
      {}[ae]+[bd]-[ad]&\mathrel{\bowtie}[ce]+[bd]-[cd],\\
      {}[cd]+[be]-[ce]&\mathrel{\bowtie}[ad]+[be]-[ae],\\
      {}[ab]+[cd]-[ad]&\mathrel{\bowtie}[ab]+[ce]-[ae],\\
      {}[cb]+[ae]-[ce]&\mathrel{\bowtie}[cb]+[ad]-[cd].
    \end{aligned}
  \]
  For example, Definition~\ref{def:game-mul} gives
  \[
    \mulopt(a,e;b,d)=ae+bd-ad,
    \qquad
    \mulopt(c,e;b,d)=ce+bd-cd.
  \]
  Hence the first rearranged inequality is exactly
  \[
    [\mulopt(a,e;b,d)]\mathrel{\bowtie}[\mulopt(c,e;b,d)],
  \]
  which is the class-level form of rule \textup{(a)}.
  The other three lines give rules \textup{(b)}--\textup{(d)} in the same way.
  The argument works for both weak and strict order because translation preserves both in
  the ordered additive group.

  We prove rule \textup{(g)},
  \[
    \operatorname{Cross}_{\GameLt}(a,b;d,e)
      \ \wedge\ bd\GameEq cd
      \ \wedge\ be\GameEq ce
      \quad\Longrightarrow\quad
    \operatorname{Cross}_{\GameLt}(a,c;d,e);
  \]
  the other strict rules \textup{(e)}, \textup{(f)}, and \textup{(h)} are analogous.
  The first hypothesis means
  \[
    ae+bd\GameLt ad+be.
  \]
  By Theorem~\ref{thm:game-quotient-ordered-add-group}, this is the strict
  quotient inequality
  \[
    [ae]+[bd]<[ad]+[be].
  \]
  The two equivalence hypotheses give the class equalities
  \[
    [bd]=[cd],
    \qquad
    [be]=[ce].
  \]
  Substituting them into the strict inequality gives
  \[
    [ae]+[cd]<[ad]+[ce].
  \]
  Reflecting strict order from the quotient yields
  \[
    ae+cd\GameLt ad+ce,
  \]
  which is exactly
  \(\operatorname{Cross}_{\GameLt}(a,c;d,e)\).
\end{proof}

\begin{lemma}[Terminal product-option bounds]
  \label{lem:mulopt-bounds}
  \lean{mulOpt4_LL_lt_product,mulOpt4_RR_lt_product,product_lt_mulOpt4_LR,product_lt_mulOpt4_RL}
  \leanok
  The following statements hold:
  \[
    \begin{array}{ll}
      (LL)&ay^L\GameEq by^L,\ C_L(a^L,b;y,y^L)
        \Longrightarrow \mulopt(a^L,y;a,y^L)\GameLt by,\\[1mm]
      (RR)&ay^R\GameEq by^R,\ C_R(b,a^R;y,y^R)
        \Longrightarrow \mulopt(a^R,y;a,y^R)\GameLt by,\\[1mm]
      (LR)&by^R\GameEq ay^R,\ C_R(b^L,a;y,y^R)
        \Longrightarrow ay\GameLt\mulopt(b^L,y;b,y^R),\\[1mm]
      (RL)&by^L\GameEq ay^L,\ C_L(a,b^R;y,y^L)
        \Longrightarrow ay\GameLt\mulopt(b^R,y;b,y^L).
    \end{array}
  \]
\end{lemma}
\begin{proof}
  \leanok
  \uses{def:game-mul,def:game-quotient,thm:game-quotient-ordered-add-group,def:cross-conditions}
  We prove only the \(LL\) statement.  Its hypotheses are
  \[
    ay^L\GameEq by^L
    \quad\text{and}\quad
    C_L(a^L,b;y,y^L):
    a^Ly+by^L\GameLt a^Ly^L+by.
  \]
  In the quotient \(\mathcal G/{\GameEq}\), add \(-[a^Ly^L]\) to both sides and use
  \([by^L]=[ay^L]\).  This gives
  \[
    [a^Ly]+[ay^L]-[a^Ly^L]<[by],
  \]
  whose left side is \([\mulopt(a^L,y;a,y^L)]\) by
  Definition~\ref{def:game-mul}.  The quotient order in
  Theorem~\ref{thm:game-quotient-ordered-add-group} therefore gives
  \[
    \mulopt(a^L,y;a,y^L)\GameLt by.
  \]
  The \(RR\), \(LR\), and \(RL\) statements follow from the same quotient
  rearrangement with their displayed equivalence and cross-inequality hypotheses.
\end{proof}

\begin{theorem}[Simultaneous Conway theorem]
  \label{thm:abc-main}
  \lean{Conway.Holds,Conway.ABC_main}
  \leanok
  The following assertions hold simultaneously for all games satisfying the
  displayed hypotheses:
  \begin{enumerate}
    \item \(A(x,y)\): if \(x,y\) are surreal, then \(xy\) is surreal;
    \item \(B(x_1,x_2,y)\): if \(x_1\GameEq x_2\) and all three games are surreal, then
      \(x_1y\GameEq x_2y\);
    \item \(C(x_1,x_2,y)\): if \(x_1\GameLt x_2\) and all three games are surreal, then
      \(C_L(x_1,x_2;y,y^L)\) holds for every \(y^L\), and
      \(C_R(x_1,x_2;y,y^R)\) holds for every \(y^R\).
  \end{enumerate}
\end{theorem}
\begin{proof}
  \leanok
  \uses{def:game-order,def:game-equivalence,def:is-surreal,def:game-mul,def:game-quotient,thm:game-quotient-ordered-add-group,def:dm-goals,def:cross-conditions,thm:game-order-calculus,thm:between-options,thm:surreal-totality,thm:surreal-add-closure,thm:game-mul-laws,thm:cross-calculus,lem:mulopt-bounds}
  Apply well-founded induction with the goal order of
  Definition~\ref{def:dm-goals}.  The induction hypothesis is available for
  every \(A\)-, \(B\)-, or \(C\)-assertion whose multiset measure is smaller.
  The comparisons in that definition justify replacing a coordinate by an
  option, replacing one birthday by two smaller birthdays, and deleting an
  unused coordinate.

  Consider \(A(x,y)\).  By Definition~\ref{def:is-surreal}, it is enough to
  prove that every option of \(xy\) is surreal and that every left option is
  strictly below every right option.  Definition~\ref{def:game-mul} gives the
  four option families.  If \(x'\) and \(y'\) are options of \(x\) and \(y\),
  respectively, the corresponding option has the form
  \[
    \mulopt(x',y;x,y')=x'y+xy'-x'y'.
  \]
  The smaller assertions \(A(x',y)\), \(A(x,y')\), and \(A(x',y')\) prove
  that its three product terms are surreal.  Closure under addition and
  negation, Theorem~\ref{thm:surreal-add-closure}, then proves the whole
  expression surreal.

  For separation put \(P(u,v)=\mulopt(u,y;x,v)\).  By
  Definition~\ref{def:game-mul}, a left option of \(xy\) is
  \[
    P(x^L,y^L)\quad(LL),\qquad P(x^R,y^R)\quad(RR),
  \]
  and a right option is
  \[
    P(x^L,y^R)\quad(LR),\qquad P(x^R,y^L)\quad(RL).
  \]
  Consider first \(P(x^L_1,y^L)\) and \(P(x^L_2,y^R)\).  Since the first is a
  left option and the second is a right option, we need to prove
  \[
    P(x^L_1,y^L)\GameLt P(x^L_2,y^R).
  \]
  Theorem~\ref{thm:surreal-totality} compares \(x^L_1\) and \(x^L_2\).
  Also, Theorem~\ref{thm:between-options} and strict transitivity from
  Theorem~\ref{thm:game-order-calculus} give
  \(y^L\GameLt y\GameLt y^R\), hence \(y^L\GameLt y^R\).
  The three possible comparisons of the two left options give
  \[
    \begin{array}{ll}
    x^L_1\GameLt x^L_2:
      &P(x^L_1,y^L)\GameLt P(x^L_2,y^L)
         \GameLt P(x^L_2,y^R),\\[1mm]
    x^L_1\GameEq x^L_2:
      &P(x^L_1,y^L)\GameLt P(x^L_1,y^R)
         \GameLe P(x^L_2,y^R),\\[1mm]
    x^L_2\GameLt x^L_1:
      &P(x^L_1,y^L)\GameLt P(x^L_1,y^R)
         \GameLt P(x^L_2,y^R).
    \end{array}
  \]
  Here is the first branch in detail.  Assume \(x^L_1\GameLt x^L_2\).  The
  smaller assertion \(C(x^L_1,x^L_2,y)\) gives
  \[
    C_L(x^L_1,x^L_2;y,y^L)
      =\operatorname{Cross}_{\GameLt}(x^L_1,x^L_2;y^L,y).
  \]
  Substitute
  \((a,c;d,e;b)=(x^L_1,x^L_2;y^L,y;x)\) in rule \textup{(a)} of
  Theorem~\ref{thm:cross-calculus}.  It becomes
  \[
    \mulopt(x^L_1,y;x,y^L)\GameLt\mulopt(x^L_2,y;x,y^L),
    \qquad\text{that is,}\qquad
    P(x^L_1,y^L)\GameLt P(x^L_2,y^L).
  \]
  For the second step, the smaller assertion \(C(y^L,y^R,x)\), evaluated at
  the left option \(x^L_2\) of \(x\), gives
  \[
    \operatorname{Cross}_{\GameLt}(y^L,y^R;x^L_2,x).
  \]
  Applying rule \textup{(e)} of Theorem~\ref{thm:cross-calculus} gives
  \(\operatorname{Cross}_{\GameLt}(x^L_2,x;y^L,y^R)\).  Rule \textup{(c)},
  now with
  \((a,c;d,e;b)=(x^L_2,x;y^L,y^R;y)\), gives
  \[
    P(x^L_2,y^L)\GameLt P(x^L_2,y^R).
  \]
  Strict transitivity therefore proves the first chain in the display.

  If \(x^L_1\GameEq x^L_2\), the same \(y\)-slot argument first gives
  \(P(x^L_1,y^L)\GameLt P(x^L_1,y^R)\).  The two smaller assertions
  \(B(x^L_1,x^L_2,y)\) and \(B(x^L_1,x^L_2,y^R)\) give
  \[
    x^L_1y\GameEq x^L_2y,
    \qquad
    x^L_1y^R\GameEq x^L_2y^R.
  \]
  Hence, in \(\Games/{\GameEq}\),
  \[
    [P(x^L_1,y^R)]
      =[x^L_1y+xy^R-x^L_1y^R]
      =[x^L_2y+xy^R-x^L_2y^R]
      =[P(x^L_2,y^R)].
  \]
  Thus \(P(x^L_1,y^R)\GameEq P(x^L_2,y^R)\), and the strict--weak
  composition law proves the middle chain.

  Finally, assume \(x^L_2\GameLt x^L_1\).  After the same strict
  \(y\)-slot movement at \(x^L_1\), the smaller assertion
  \(C(x^L_2,x^L_1,y)\) gives
  \[
    C_R(x^L_2,x^L_1;y,y^R)
      =\operatorname{Cross}_{\GameLt}(x^L_2,x^L_1;y,y^R).
  \]
  Rule \textup{(b)} of Theorem~\ref{thm:cross-calculus},
  with \((a,c;d,e;b)=(x^L_2,x^L_1;y,y^R;x)\), reverses the \(x\)-slot
  direction and gives
  \[
    P(x^L_1,y^R)\GameLt P(x^L_2,y^R).
  \]
  Strict transitivity proves the last chain.  All the \(B\)- and
  \(C\)-assertions just invoked are smaller by the one-to-two birthday
  replacement in Definition~\ref{def:dm-goals}.

  The \(LL\)-versus-\(RL\), \(RR\)-versus-\(LR\), and
  \(RR\)-versus-\(RL\) comparisons are proved analogously.  This proves
  \(A(x,y)\).

  Now consider \(B(a,b,y)\).  Assume \(a\GameEq b\).  By
  Definition~\ref{def:game-order}, to prove \(ay\GameLe by\) one must rule out
  \(by\GameLe\lambda\) for every left option \(\lambda\) of \(ay\), and rule
  out \(\rho\GameLe ay\) for every right option \(\rho\) of \(by\).

  Definition~\ref{def:game-mul} gives two kinds of left option of \(ay\):
  \[
    \begin{aligned}
      \lambda_{LL}&=\mulopt(a^L,y;a,y^L)
        =a^Ly+ay^L-a^Ly^L,\\
      \lambda_{RR}&=\mulopt(a^R,y;a,y^R)
        =a^Ry+ay^R-a^Ry^R.
    \end{aligned}
  \]
  We must prove \(\lambda_{LL}\GameLt by\) and
  \(\lambda_{RR}\GameLt by\); either strict inequality contains the
  corresponding conclusion \(\neg(by\GameLe\lambda)\).

  For \(\lambda_{LL}\), the smaller assertion \(B(a,b,y^L)\) gives
  \[
    ay^L\GameEq by^L.
  \]
  The comparison \(a^L\GameLt a\), together with \(a\GameLe b\) from
  \(a\GameEq b\), gives \(a^L\GameLt b\).  Hence the smaller assertion
  \(C(a^L,b,y)\) gives the explicit cross inequality
  \[
    C_L(a^L,b;y,y^L):\qquad
    a^Ly+by^L\GameLt a^Ly^L+by.
  \]
  These are exactly the two hypotheses of the \(LL\) implication in
  Lemma~\ref{lem:mulopt-bounds}, so
  \[
    \lambda_{LL}=a^Ly+ay^L-a^Ly^L\GameLt by.
  \]

  For \(\lambda_{RR}\), the smaller assertion \(B(a,b,y^R)\) gives
  \[
    ay^R\GameEq by^R.
  \]
  Now \(b\GameLe a\) follows from \(a\GameEq b\), while
  \(a\GameLt a^R\); hence \(b\GameLt a^R\).  The smaller assertion
  \(C(b,a^R,y)\) therefore gives
  \[
    C_R(b,a^R;y,y^R):\qquad
    by^R+a^Ry\GameLt by+a^Ry^R.
  \]
  The \(RR\) implication in Lemma~\ref{lem:mulopt-bounds} now yields
  \[
    \lambda_{RR}=a^Ry+ay^R-a^Ry^R\GameLt by.
  \]
  Every assertion used here is smaller by Definition~\ref{def:dm-goals} and
  Lemma~\ref{lem:birthday-options}.

  The two right-option cases are analogous.  Explicitly, for the two right
  options of \(by\), one must prove
  \[
    \begin{aligned}
      ay&\GameLt \rho_{LR}
        =\mulopt(b^L,y;b,y^R)=b^Ly+by^R-b^Ly^R,\\
      ay&\GameLt \rho_{RL}
        =\mulopt(b^R,y;b,y^L)=b^Ry+by^L-b^Ry^L.
    \end{aligned}
  \]
  These follow analogously from the \(LR\) and \(RL\) implications of
  Lemma~\ref{lem:mulopt-bounds}.  The four exclusions prove \(ay\GameLe by\).

  Interchanging \(a\) and \(b\) gives \(by\GameLe ay\).  This does not invoke
  the induction hypothesis at the unchanged measure: the equality
  \[
    \mu_B(a,b,y)=\mu_B(b,a,y)
  \]
  from Definition~\ref{def:dm-goals} identifies the two collections of
  strictly smaller assertions, and symmetry of \(a\GameEq b\) permits the
  same one-sided proof to be reused.  The two weak inequalities are exactly
  \(ay\GameEq by\) by Definition~\ref{def:game-equivalence}.

  Finally consider \(C(x_1,x_2,y)\), with \(x_1\GameLt x_2\).  The assertion
  to be proved is that, for every left option \(y^L\) and right option \(y^R\)
  of \(y\),
  \[
    \begin{aligned}
      C_L(x_1,x_2;y,y^L):\quad
        &x_1y+x_2y^L\GameLt x_1y^L+x_2y,\\
      C_R(x_1,x_2;y,y^R):\quad
        &x_1y^R+x_2y\GameLt x_1y+x_2y^R.
    \end{aligned}
  \]
  Unfolding Definition~\ref{def:game-order} gives the following bridge
  alternative: either a right option \(x_1^R\) satisfies
  \(x_1^R\GameLe x_2\), or a left option \(x_2^L\) satisfies
  \(x_1\GameLe x_2^L\).  Indeed, if neither existed, the two universal clauses
  in that definition would give \(x_2\GameLe x_1\), contradicting
  \(x_1\GameLt x_2\).

  Suppose first that \(x_1^R\GameLe x_2\).  The smaller assertion
  \(A(x_1,y)\) makes \(x_1y\) surreal.  Its product options are those of
  Definition~\ref{def:game-mul}.  For each left option \(y^L\) of \(y\), the
  \(RL\) product option lies above \(x_1y\), so
  \[
    x_1y\GameLt x_1^Ry+x_1y^L-x_1^Ry^L.
  \]
  Work in the quotient of Definition~\ref{def:game-quotient}, whose ordered additive group
  structure is given by Theorem~\ref{thm:game-quotient-ordered-add-group}.  Adding
  \([x_1^Ry^L]\) gives
  \[
    [x_1y]+[x_1^Ry^L]<[x_1^Ry]+[x_1y^L].
  \]
  Reflection of strict order from the quotient gives
  \[
    \begin{aligned}
      C_L(x_1,x_1^R;y,y^L)
        &=\operatorname{Cross}_{\GameLt}(x_1,x_1^R;y^L,y),\\
        &\text{i.e.}\quad
          x_1y+x_1^Ry^L\GameLt x_1y^L+x_1^Ry.
    \end{aligned}
  \]
  Similarly, for a right option \(y^R\),
  the \(RR\) product option lies below \(x_1y\):
  \[
    x_1^Ry+x_1y^R-x_1^Ry^R\GameLt x_1y.
  \]
  Adding \([x_1^Ry^R]\) and reflecting strict order gives
  \[
    \begin{aligned}
      C_R(x_1,x_1^R;y,y^R)
        &=\operatorname{Cross}_{\GameLt}(x_1,x_1^R;y,y^R),\\
        &\text{i.e.}\quad
          x_1y^R+x_1^Ry\GameLt x_1y+x_1^Ry^R.
    \end{aligned}
  \]
  These are called the adjacent conditions because their two endpoints are
  \(x_1\) and its option \(x_1^R\).  The required conditions displayed above
  have endpoints \(x_1,x_2\), so it remains to move the second endpoint from
  \(x_1^R\) to \(x_2\).

  By Theorem~\ref{thm:surreal-totality}, either
  \(x_1^R\GameLt x_2\), \(x_1^R\GameEq x_2\), or
  \(x_2\GameLt x_1^R\).  If \(x_1^R\GameLt x_2\), the smaller assertion
  \(C(x_1^R,x_2,y)\) supplies
  \[
    \begin{aligned}
      C_L(x_1^R,x_2;y,y^L)
        &=\operatorname{Cross}_{\GameLt}(x_1^R,x_2;y^L,y),\\
        &\text{i.e.}\quad
          x_1^Ry+x_2y^L\GameLt x_1^Ry^L+x_2y,\\[1mm]
      C_R(x_1^R,x_2;y,y^R)
        &=\operatorname{Cross}_{\GameLt}(x_1^R,x_2;y,y^R),\\
        &\text{i.e.}\quad
          x_1^Ry^R+x_2y\GameLt x_1^Ry+x_2y^R.
    \end{aligned}
  \]
  Rule \textup{(f)} of Theorem~\ref{thm:cross-calculus} now composes
  \[
    \begin{gathered}
      \operatorname{Cross}_{\GameLt}(x_1,x_1^R;y^L,y)
      \quad\text{and}\quad
      \operatorname{Cross}_{\GameLt}(x_1^R,x_2;y^L,y),\\
      \operatorname{Cross}_{\GameLt}(x_1,x_1^R;y,y^R)
      \quad\text{and}\quad
      \operatorname{Cross}_{\GameLt}(x_1^R,x_2;y,y^R).
    \end{gathered}
  \]
  The results are precisely
  \(C_L(x_1,x_2;y,y^L)\) and \(C_R(x_1,x_2;y,y^R)\).

  If \(x_1^R\GameEq x_2\), the smaller \(B\)-assertions give
  \[
    x_1^Ry\GameEq x_2y,\qquad
    x_1^Ry^L\GameEq x_2y^L,\qquad
    x_1^Ry^R\GameEq x_2y^R.
  \]
  Replacing these three products in the two adjacent inequalities gives,
  respectively,
  \[
    x_1y+x_2y^L\GameLt x_1y^L+x_2y,
    \qquad
    x_1y^R+x_2y\GameLt x_1y+x_2y^R,
  \]
  which are again the two required conditions.  The remaining alternative
  \(x_2\GameLt x_1^R\) contradicts \(x_1^R\GameLe x_2\).

  The alternative \(x_1\GameLe x_2^L\) is analogous, using the adjacent
  conditions for \(x_2^L,x_2\) and then moving the first endpoint from
  \(x_2^L\) to \(x_1\).
\end{proof}

\begin{theorem}[Conway A: closure of the product]
  \label{thm:conway-a}
  \lean{Conway.conway_A}
  \leanok
  If \(x\) and \(y\) are surreal games, then \(xy\) is surreal.
\end{theorem}
\begin{proof}
  \leanok
  \uses{thm:abc-main}
  Apply the \(A(x,y)\) clause of Theorem~\ref{thm:abc-main} to the two
  surrealness hypotheses.  Its conclusion is precisely that \(xy\) is
  surreal.
\end{proof}

\begin{theorem}[Conway B: product respects equivalence]
  \label{thm:conway-b}
  \lean{Conway.conway_B}
  \leanok
  If \(x_1,x_2,y\) are surreal and \(x_1\GameEq x_2\), then
  \(x_1y\GameEq x_2y\).
\end{theorem}
\begin{proof}
  \leanok
  \uses{thm:abc-main}
  Apply the \(B(x_1,x_2,y)\) clause of Theorem~\ref{thm:abc-main}, with the
  three surrealness hypotheses and \(x_1\GameEq x_2\).  The conclusion is
  \(x_1y\GameEq x_2y\).
\end{proof}

\begin{theorem}[Conway C: strict cross inequalities]
  \label{thm:conway-c}
  \lean{Conway.conway_C}
  \leanok
  If \(x_1,x_2,y\) are surreal and \(x_1\GameLt x_2\), then all the conditions
  \(C_L(x_1,x_2;y,y^L)\) and \(C_R(x_1,x_2;y,y^R)\) hold.
\end{theorem}
\begin{proof}
  \leanok
  \uses{thm:abc-main}
  Apply the \(C(x_1,x_2,y)\) clause of Theorem~\ref{thm:abc-main}, with the
  three surrealness hypotheses and \(x_1\GameLt x_2\).  Its two conclusions
  are precisely the required left and right cross inequalities.
\end{proof}

\section{Congruence and associativity}

\begin{theorem}[Two-variable multiplicative congruence]
  \label{thm:game-mul-congr}
  \lean{Game.mul_congr}
  \leanok
  If \(x_1\GameEq x_2\) and \(y_1\GameEq y_2\), and all four games are surreal, then
  \[
    x_1y_1\GameEq x_2y_2.
  \]
\end{theorem}
\begin{proof}
  \leanok
  \uses{thm:game-order-calculus,thm:game-mul-laws,thm:conway-b}
  Conway \(B\), Theorem~\ref{thm:conway-b}, applied to
  \(x_1\GameEq x_2\) gives
  \(x_1y_1\GameEq x_2y_1\).  To replace the second factor, use product
  commutativity from Theorem~\ref{thm:game-mul-laws}, apply
  Theorem~\ref{thm:conway-b} to \(y_1\GameEq y_2\) with common factor \(x_2\),
  and commute back.  This gives
  \(x_2y_1\GameEq x_2y_2\).  Transitivity of game equivalence in
  Theorem~\ref{thm:game-order-calculus} composes the two replacements.
\end{proof}

\begin{theorem}[Associativity on surreal games]
  \label{thm:game-mul-assoc}
  \lean{Game.mul_assoc_of_isSurreal}
  \leanok
  If \(a,b,c\) are surreal games, then
  \[
    (ab)c\GameEq a(bc).
  \]
\end{theorem}
\begin{proof}
  \leanok
  \uses{def:game-inductions,lem:birthday-options,def:game-mul,def:game-neg,def:game-quotient,thm:game-quotient-ordered-add-group,def:is-surreal,thm:game-order-calculus,thm:option-extensionality,thm:game-neg-order,thm:game-mul-laws,thm:game-distrib,lem:game-mul-neg}
  Use the three-game well-founded induction of
  Definition~\ref{def:game-inductions}.  Lemma~\ref{lem:birthday-options}
  makes every option replacement smaller.  We strengthen the induction
  statement as follows: associativity may be used for any triple whose three
  birthdays are bounded by those of \((a,b,c)\), provided at least one bound
  is strict.  Thus, if \(a',b',c'\) are chosen options, the induction
  hypothesis contains all seven equivalences
  \[
    \begin{gathered}
      (a'b)c\GameEq a'(bc),\qquad
      (ab')c\GameEq a(b'c),\qquad
      (ab)c'\GameEq a(bc'),\\
      (a'b')c\GameEq a'(b'c),\qquad
      (a'b)c'\GameEq a'(bc'),\qquad
      (ab')c'\GameEq a(b'c'),\\
      (a'b')c'\GameEq a'(b'c').
    \end{gathered}
  \]

  We next compare a nested product option under the two bracketings.  For
  options \(a',b',c'\), we claim that:
  \[
    \mulopt\bigl(\mulopt(a',b;a,b'),c;ab,c'\bigr)
      \GameEq
    \mulopt\bigl(a',bc;a,\mulopt(b',c;b,c')\bigr).
  \]
  We prove this by expanding the left side and the right side separately.
  First, \(\mulopt(a',b;a,b')=a'b+ab'-a'b'\), so distributivity gives
  \[
    \begin{aligned}
      &\left[
        \mulopt\bigl(\mulopt(a',b;a,b'),c;ab,c'\bigr)
      \right]\\
      &\quad=
        \bigl[\mulopt(a',b;a,b')c\bigr]+[(ab)c']
          -\bigl[\mulopt(a',b;a,b')c'\bigr]\\
        &=[(a'b)c]+[(ab')c]-[(a'b')c]+[(ab)c']\\
        &\quad-[(a'b)c']-[(ab')c']+[(a'b')c'].
    \end{aligned}
  \]
  This is where the second equivalence of Lemma~\ref{lem:game-mul-neg} is
  used: right-distributivity produces the terms
  \[
    [(-(a'b'))c]
      =[-((a'b')c)],
    \qquad
    [(-(a'b'))c']
      =[-((a'b')c')].
  \]
  Substituting the seven strengthened induction hypotheses into this equality
  gives
  \[
    \begin{aligned}
      &\left[
        \mulopt\bigl(\mulopt(a',b;a,b'),c;ab,c'\bigr)
      \right]\\
        &\quad=[a'(bc)]+[a(b'c)]-[a'(b'c)]+[a(bc')]\\
        &\quad-[a'(bc')]-[a(b'c')]+[a'(b'c')].
    \end{aligned}
  \]

  For the other bracketing,
  \(\mulopt(b',c;b,c')=b'c+bc'-b'c'\).  Expanding directly gives
  \[
    \begin{aligned}
      &\left[
        \mulopt\bigl(a',bc;a,\mulopt(b',c;b,c')\bigr)
      \right]\\
      &\quad=[a'(bc)]
        +\bigl[a\mulopt(b',c;b,c')\bigr]
        -\bigl[a'\mulopt(b',c;b,c')\bigr]\\
        &=[a'(bc)]+[a(b'c)]+[a(bc')]-[a(b'c')]\\
        &\quad-[a'(b'c)]-[a'(bc')]+[a'(b'c')].
    \end{aligned}
  \]
  Here the first equivalence of Lemma~\ref{lem:game-mul-neg} is used in the
  two terms
  \[
    [a(-(b'c'))]=[-a(b'c')],
    \qquad
    [a'(-(b'c'))]=[-a'(b'c')].
  \]
  The right sides of the last two class expansions differ only by the order
  and association of their summands.  The additive commutative group structure
  of Theorem~\ref{thm:game-quotient-ordered-add-group} therefore gives
  \[
    \mulopt\bigl(\mulopt(a',b;a,b'),c;ab,c'\bigr)
      \GameEq
    \mulopt\bigl(a',bc;a,\mulopt(b',c;b,c')\bigr).
  \]
  This proves the claim.

  We now apply this claim to the option families of
  Definition~\ref{def:game-mul}.  First choose an \(LL\) option of \((ab)c\)
  whose selected left option of \(ab\) is itself \(LL\).  For
  \(a^L\in A^L,b^L\in B^L,c^L\in C^L\), it is
  \[
    L=\mulopt\bigl(\mulopt(a^L,b;a,b^L),c;ab,c^L\bigr).
  \]
  The game \(\mulopt(b^L,c;b,c^L)\) is an \(LL\) left option of \(bc\), so
  \[
    L'=\mulopt\bigl(a^L,bc;a,\mulopt(b^L,c;b,c^L)\bigr)
  \]
  is an \(LL\) left option of \(a(bc)\).  The preceding calculation with
  \((a',b',c')=(a^L,b^L,c^L)\) gives \(L\GameEq L'\).

  For a mixed example, choose an \(LR\) right option of \((ab)c\), again
  using an \(LL\) left option of \(ab\).  With \(c^R\in C^R\), it is
  \[
    R=\mulopt\bigl(\mulopt(a^L,b;a,b^L),c;ab,c^R\bigr).
  \]
  Now \(\mulopt(b^L,c;b,c^R)\) is an \(LR\) right option of \(bc\), and
  \[
    R'=\mulopt\bigl(a^L,bc;a,\mulopt(b^L,c;b,c^R)\bigr)
  \]
  is the corresponding \(LR\) right option of \(a(bc)\).  The calculation
  with \((a',b',c')=(a^L,b^L,c^R)\) gives \(R\GameEq R'\).

  The remaining cases choose \(a^R,b^R,c^L,c^R\) according to the outer and
  inner \(LL,RR,LR,RL\) families.  Reversing the calculation matches options
  from \(a(bc)\) back to \((ab)c\).  Definition~\ref{def:is-surreal} shows
  that every chosen option of \(a,b,c\) is surreal, so the seven strengthened
  recursive associativity statements apply.  Thus the left and right option
  lists match up to equivalence in both directions.
  Theorem~\ref{thm:option-extensionality} yields
  \((ab)c\GameEq a(bc)\).
\end{proof}
